%=========================================================================
%  THEME 2 — "SIDEBAR"
%
%  Zweispaltig: getönte Seitenspalte über die volle Seitenhöhe, grosses
%  Portrait, Schlagwort-Chips. Warme Sandfläche statt schwarzem Balken,
%  Source Sans Pro als etwas schmalere, laufruhige Textschrift.
%  Für HR-Screening gedacht: alles Faktische steht links, die Nachweise
%  rechts.
%=========================================================================

\usepackage[left=1.3cm, right=1.3cm, top=1.15cm, bottom=1.25cm,
            headheight=13pt, headsep=10pt, footskip=0.8cm]{geometry}

%=========================================================================
%  GEMEINSAME BASIS — Pakete, Icons, Hilfsmakros.
%  Enthält bewusst KEINE Gestaltung. Alles Gestalterische (Farben, Schrift,
%  Raster, Abstände) liegt im jeweiligen Theme unter templates/themes/.
%=========================================================================

\usepackage[british, ngerman]{babel}
\usepackage{etoolbox}
\usepackage{xcolor}
\usepackage{fontspec}
\usepackage[final]{microtype}

%-------------------------------------------------------------------------
%  Echte Leerzeichen im PDF. TeX setzt Wortzwischenräume sonst
%  nur als Abstand. Ältere PDF-Leser (z.B. pdf.js 2.x im PDF-Viewer von
%  VS Code) und manche Bewerbungsportale lesen den Text dann ohne Leerzeichen
%  aus ("Bei500TaktenproMinute"). Mit 2 schreibt XeTeX die Leerzeichen als
%  Glyphen mit. Seitenzahlen und Blocksatz bleiben gleich, einzelne Absätze
%  brechen etwas anders um (weniger Trennungen).
%  Gemessen mit pdf.js 2.6, Anteil einzeln lesbarer Wörter: Briefe vorher
%  3-5 %, nachher 97-99 %. Im Lebenslauf stieg der Anteil ebenfalls deutlich.
%  Rest: Zeilen in der Kompetenzentabelle und einzelne Zeilen mit Bindestrich
%  verlieren die Zeichen weiterhin (XeTeX-intern, nicht weiter eingegrenzt).
%  pdftotext, pdf.js 4 und gängige Parser setzen dort ohnehin Leerzeichen aus
%  den Abständen, ohne Doppelungen.
%  ACHTUNG, XeTeX 0.999993 stürzt damit ab (Segmentation fault), wenn
%    - irgendwo \spaceskip ungleich null ist (deshalb der Umbau bei ragged2e
%      unten) oder
%    - \RaggedRight/\bodyflow mitten in einem schon begonnenen Absatz
%      eingeschaltet wird ("Text {\bodyflow ...}"). Immer am Absatzanfang.
%  Und: Farbwechsel mitten im Fliesstext in derselben Schrift verrutschen
%  (das farbige Wort bleibt schwarz, der Text danach wird farbig). Abhilfe:
%  \kern0pt an beiden Enden des farbigen Stücks, z.B.
%  {\color{red}\kern0pt Text\kern0pt} (so \platz in feedback.tex.
%  Farbe zusammen mit Schriftwechsel (Ortsangaben, Symbole) ist nicht betroffen.
%  Stürzt ein Build nach einer Vorlagenänderung ab, zuerst diese beiden Punkte
%  prüfen. Notausgang: hier \XeTeXinterwordspaceshaping=0 setzen.
%-------------------------------------------------------------------------
\XeTeXinterwordspaceshaping=2
% Als Zeichen geschrieben wird nur ein normaler Wortabstand. Nach Grossbuch-
% staben (sfcode 999) und Satzzeichen (ohne \frenchspacing) setzt TeX einen
% leicht abweichenden Abstand und XeTeX lässt das Zeichen weg ("SAPHANACloud").
% Deshalb alle auf 1000 -- im Deutschen ohnehin üblich, im Satzbild unsichtbar.
\count255=`A
\loop\sfcode\count255=1000 \ifnum\count255<`Z \advance\count255 by 1 \repeat
\sfcode`Ä=1000 \sfcode`Ö=1000 \sfcode`Ü=1000
\AtBeginDocument{\frenchspacing}
% Unsichtbares Leerzeichen ohne Breite: dort, wo im Satz kein Wortabstand
% steht und Leser trotzdem einen brauchen -- zwischen zwei Spalten
% ("EntwicklungPython"), nach Aufzählungszeichen und am Schriftwechsel nach
% den Leads ("Prüfstandsdatenplattform:Von"). Ändert das Layout nicht.
\newcommand{\cvrealspace}{\ifnum\XeTeXcharglyph"20>0
    \rlap{\XeTeXglyph\XeTeXcharglyph"20\relax}\fi}
\usepackage{tikz}
\usetikzlibrary{calc}
\usepackage{graphicx}
\usepackage{enumitem}
\usepackage{fancyhdr}
\usepackage{tabularx}
\usepackage{array}
\usepackage[fixed]{fontawesome5}
\usepackage{paracol}
\usepackage{setspace}
\usepackage{ragged2e}
% ragged2e hält im Flattersatz die Wortabstände starr, indem es bei jedem
% Schriftwechsel \spaceskip auf den natürlichen Wortabstand setzt. Zusammen
% mit echten Leerzeichen (\XeTeXinterwordspaceshaping=2, oben) stürzt XeTeX
% 0.999993 dabei ab. Gleicher Effekt ohne \spaceskip:
% Dehnung und Stauchung des Wortabstands in der Schrift selbst auf null setzen
% (\fontdimen3/4). Der Wortabstand ist damit derselbe wie bei ragged2e.
% Nebenwirkung: Die Änderung gilt für die ganze Satzlauf-Instanz der Schrift.
% Unschädlich, weil die Briefe im Blocksatz vor dem Lebenslauf gesetzt werden
% bzw. in eigenen Läufen (make dossier).
\makeatletter
\renewcommand{\@raggedtwoe@everyselectfont}{%
    \spaceskip\z@
    \if@raggedtwoe@spaceskip
        \fontdimen3\font=\z@
        \fontdimen4\font=\z@
    \fi}
\makeatother
\usepackage{needspace}
\usepackage{accsupp}

% Textschicht für Bewerbungsportale. SmartRecruiters & Co.
% lesen den Lebenslauf aus dem PDF aus. Ohne diese Zeile kamen gesperrte
% Überschriften als "BER UFS E R FAH R UNG" und jeder Bindestrich als
% geschützter Bindestrich (U+2011) an. \XeTeXgenerateactualtext legt zu
% jedem Textstück den Originaltext ab; am Druckbild ändert sich nichts.
\ifdefined\XeTeXgenerateactualtext\XeTeXgenerateactualtext=1\fi


%-------------------------------------------------------------------------
%  Sprache
%
%  \cvlanguage ist ein babel-Sprachname und wird VOR dem Theme gesetzt --
%  in applications/<app>/meta.tex, denn die Sprache gehoert zur Ziel-
%  bewerbung und nicht zur Gestaltung. Voreinstellung ngerman: alle
%  bestehenden Bewerbungen bauen unveraendert, ohne dass ihre meta.tex
%  etwas davon wissen muss. Die babel-Optionen oben laden alle drei
%  Sprachen; die letzte ist die Hauptsprache, \selectlanguage schaltet um.
%  Fuer Englisch "british" verwenden, nicht "english": babel warnt vor der
%  mehrdeutigen Option, und british liefert das europaeische Datumsformat.
%-------------------------------------------------------------------------

\providecommand{\cvlanguage}{ngerman}

%-------------------------------------------------------------------------
%  Beschriftungen der Abschnitte -- Voreinstellung Deutsch.
%  Eine uebersetzte Inhaltsdatei (templates/themes/cv-content-en.tex)
%  setzt sie per \renewcommand um; kein Theme enthaelt mehr festen Text.
%-------------------------------------------------------------------------

%  Briefdatum. Voreinstellung ist babels \today. Die englische Inhalts-
%  datei setzt es um: babels british schreibt "16th September 2026", und
%  die Ordnungszahl ist im Geschaeftsbrief unueblich.
\providecommand{\lettertoday}{\today}

\providecommand{\LblSkills}{Kompetenzen}
\providecommand{\LblExperience}{Berufserfahrung}
\providecommand{\LblEducation}{Ausbildung}
\providecommand{\LblFurther}{Weiterbildung}
\providecommand{\LblLanguages}{Sprachen}
\providecommand{\LblContact}{Kontakt}
\providecommand{\LblProfile}{Profil}
\providecommand{\LblCV}{Lebenslauf}
\providecommand{\LblPage}{Seite}

\setlength{\parindent}{0pt}
\setlength{\parskip}{0pt}
\emergencystretch=2em

%-------------------------------------------------------------------------
%  Typografische Hygiene
%-------------------------------------------------------------------------

% Trennung in Überschriften/Titeln unterdrücken — verhindert Unfälle wie
% "MUSTERMANN TECH-|NIK AG" quer über den Zeilenumbruch.
\newcommand{\nohyph}{\hyphenpenalty=10000\exhyphenpenalty=10000\hbadness=10000}

% Eigennamen, die nie getrennt werden sollen. \hyphenation traegt die
% Ausnahmen nur in die gerade aktive Sprache ein -- deshalb als Makro und
% nach \selectlanguage noch einmal aufgerufen.
\newcommand{\cvhyphenexceptions}{%
    % Eigene Firmen-, Orts- und Produktnamen hier ergänzen. Ohne Bindestrich
    % = nie trennen, mit Bindestrich = nur an dieser Fuge trennen.
    \hyphenation{ANSYS FastAPI Team-center Work-bench}}
\cvhyphenexceptions
% \frenchspacing NACH \selectlanguage: babel setzt den Satzabstand je
% Sprache, und british stellt den weiten englischen Abstand nach dem Punkt
% ein. Der faellt im Brief als Loecher auf und passt nicht zu den deutschen
% Dokumenten, die ngerman ohnehin auf frenchspacing stellt -- fuer Deutsch
% ist die Zeile daher wirkungslos.
\AtBeginDocument{\expandafter\selectlanguage\expandafter{\cvlanguage}%
    \cvhyphenexceptions\frenchspacing}

% Flatterrand statt Blocksatz: in schmalen Spalten erzeugt Blocksatz
% Trennungskaskaden und Löcher. \RaggedRight erlaubt Trennung noch,
% hält den Grauwert aber ruhig. \cvbodyhyphenpenalty kann die Inhaltsdatei
% überschreiben: 10000 schaltet die Trennung ab  --
% getrennte Wörter kamen in Bewerbungsportalen als "Haushaltsvoll- zugs-
% programm" an, weil die Portale den Text aus dem PDF auslesen.
\providecommand{\cvbodyhyphenpenalty}{200}
% Ohne Trennung dehnt \raggedright statt \RaggedRight: ragged2e lässt eine
% Zeile vor einem langen Wort lieber überlaufen ("Prozessdokumentation"
% ragte 13,6 pt in den Rand).
\newcommand{\bodyflow}{\RaggedRight\hyphenpenalty=\cvbodyhyphenpenalty\relax
    \ifnum\cvbodyhyphenpenalty<10000 \else\raggedright\fi}

%-------------------------------------------------------------------------
%  Icons
%-------------------------------------------------------------------------

% Die Symbolschrift hat keine Unicode-Zuordnung; ausgelesen wurde jedes
% Symbol zu einem Störzeichen vor der Adresse oder Telefonnummer ("Ȱ", "Л").
% \icotext gibt ihnen einen leeren Ersatztext -- Portale sehen nur die Angabe.
\newcommand{\icotext}[1]{\BeginAccSupp{method=plain,ActualText={}}#1\EndAccSupp{}}
\newcommand{\icoloc}{\icotext{\faMapMarker*}}
\newcommand{\icophone}{\icotext{\faPhone*}}
\newcommand{\icomail}{\icotext{\faEnvelope}}
\newcommand{\icolinked}{\icotext{\faLinkedin}}
\newcommand{\icocal}{\icotext{\faBirthdayCake}}
\newcommand{\icoheart}{\icotext{\faHeart}}
\newcommand{\icocar}{\icotext{\faCar}}

%-------------------------------------------------------------------------
%  Sechster Kontaktplatz — z.B. LinkedIn, GitHub oder der Führerschein.
%  Die Inhaltsdatei setzt Symbol, Text und optional eine URL; ist die URL
%  leer, steht der Text ohne Verlinkung da.
%-------------------------------------------------------------------------
\providecommand{\CVlinkicon}{\icolinked}
\providecommand{\CVlinktext}{}
\providecommand{\CVlinkurl}{}
\newcommand{\cvlink}{\ifdefempty{\CVlinkurl}{\CVlinktext}{\href{\CVlinkurl}{\CVlinktext}}}

%-------------------------------------------------------------------------
%  Haken für Abschnitte, die nur eine Person braucht (z.B. Referenzen).
%  Wird am Ende des Lebenslaufs ausgeführt; die Inhaltsdatei füllt ihn.
%-------------------------------------------------------------------------
\providecommand{\cvextrasections}{}

%-------------------------------------------------------------------------
%  Master-Lebenslauf und optionale Bausteine.
%
%  Die Inhaltsdatei hält ALLE Angaben. Was im Dossier erscheint, steuert
%  die meta.tex der Bewerbung:
%      \cvshow{schluessel}   zeigt einen Baustein, der sonst wegbleibt
%      \cvhide{schluessel}   lässt einen Baustein weg, der sonst erscheint
%  Die Schlüssel stehen im Kopf von themes/cv-content.tex. application.tex
%  und preview.tex definieren \cvshow und \cvhide, BEVOR sie die meta.tex
%  lesen; hier stehen sie nur als Rückfallebene für andere Treiber.
%
%  In der Inhaltsdatei:
%      \cvopt{schluessel}{...}        nur mit \cvshow oder im Master
%      \cvstd{schluessel}{...}        immer, ausser mit \cvhide (im Master immer)
%      \cvifopt{schluessel}{dann}{sonst}
%  \def\cvmaster{} (make cv-master über templates/cv-master.tex) zeigt
%  alles: Das ist der Master-Lebenslauf, aus dem je Stelle reduziert wird.
%-------------------------------------------------------------------------
\providecommand{\cvshow}[1]{\expandafter\def\csname cvshow@#1\endcsname{}}
\providecommand{\cvhide}[1]{\expandafter\def\csname cvhide@#1\endcsname{}}
\newcommand{\cvifopt}[3]{\ifdef{\cvmaster}{#2}{\ifcsdef{cvshow@#1}{#2}{#3}}}
\newcommand{\cvopt}[2]{\cvifopt{#1}{#2}{}}
\newcommand{\cvstd}[2]{\ifdef{\cvmaster}{#2}{\ifcsdef{cvhide@#1}{}{#2}}}

%-------------------------------------------------------------------------
%  Achter Kontaktplatz: GitHub (Schlüssel github in der
%  Inhaltsdatei, gedacht für Software- und Datenstellen). Leer gelassen
%  erscheint er in keinem Theme.
%-------------------------------------------------------------------------
\providecommand{\CVgithubtext}{}
\providecommand{\CVgithuburl}{}
\newcommand{\icogithub}{\icotext{\faGithub}}
\newcommand{\cvgithub}{\ifdefempty{\CVgithuburl}{\CVgithubtext}{\href{\CVgithuburl}{\CVgithubtext}}}

%-------------------------------------------------------------------------
%  Haken für einen Abschnitt direkt nach der Berufserfahrung (z.B. "Open Source & Projekte"). Leer = kein Abschnitt. Jedes
%  Theme ruft ihn in \makecvbody auf.
%-------------------------------------------------------------------------
\providecommand{\cvprojectsection}{}

%-------------------------------------------------------------------------
%  Zwischenüberschrift innerhalb einer Station (z.B. "Entwicklung" und
%  "Versuch & Prüfung"). Steht zwischen
%  zwei cvbullets-Listen und bleibt mit dem ersten Punkt darunter zusammen.
%  Vorgabe für alle Themes; kompakt setzt eine eigene Fassung.
%-------------------------------------------------------------------------
\providecommand{\cvsubhead}[1]{%
    \par\Needspace*{5\baselineskip}\addvspace{4pt}%
    {\nohyph\mediumfont #1}\par\nopagebreak
}

%-------------------------------------------------------------------------
%  Siebter Kontaktplatz — Nationalität und Bewilligung. In der Schweiz
%  gehört die Angabe bei ausländischen Bewerbenden in die Personalien
%  (Aufenthaltsbewilligung B/C, für Liechtenstein die Grenzgänger-
%  bewilligung). Leer gelassen erscheint die Zeile nicht. 
%-------------------------------------------------------------------------
\providecommand{\CVnationality}{}
\newcommand{\icoflag}{\icotext{\faPassport}}

%-------------------------------------------------------------------------
%  Neunter Kontaktplatz — Bewilligung als eigener Platz:
%  "Deutsch, Niederlassungsbewilligung C" in einem Platz lief im Kompakt-
%  Theme bis in den GitHub-Platz daneben. Mit \CVpermit stehen Nationalität
%  und Bewilligung getrennt. Leer gelassen bleibt alles wie bisher.
%-------------------------------------------------------------------------
\providecommand{\CVpermit}{}
\newcommand{\icopermit}{\icotext{\faIdCard}}

%-------------------------------------------------------------------------
%  Schrift des Brieftexts. Voreinstellung: die \normalsize des Themes,
%  also dieselbe Grösse wie der Lebenslauf. Die Inhaltsdatei kann den
%  Brief grösser setzen — bei Briefen von 200 bis 300 Wörtern bleibt die
%  Seite in 9,4 pt sonst zur Hälfte leer . Wird am Ende
%  von \makeletterhead ausgeführt, nach dem \parskip des Themes.
%-------------------------------------------------------------------------
\providecommand{\letterbodyfont}{}

%-------------------------------------------------------------------------
%  Unterschrift als Bild. Leer = keine. Die Briefe rufen \lettersignature
%  zwischen Grussformel und Namen auf; ohne Bild bleibt nur der Abstand.
%-------------------------------------------------------------------------
\providecommand{\CVsignature}{}
\newcommand{\lettersignature}{%
    \ifdefempty{\CVsignature}{\vspace{0.3cm}}{%
        \par\vspace{2pt}\noindent\includegraphics[height=0.95cm]{\CVsignature}\par\vspace{1pt}}}

%-------------------------------------------------------------------------
%  Bild-Helfer
%-------------------------------------------------------------------------
% Quadratischer Ausschnitt mit abgerundeten Ecken.
%   #1 Kantenlänge des Rahmens   #2 Eckenradius
% Zwei Stellschrauben kommen aus der Inhaltsdatei, weil sie zum Foto
% gehören und nicht zum Theme — beide relativ zur Kantenlänge, damit
% derselbe Wert in allen Themes stimmt, egal wie gross der Rahmen ist:
%   \CVphotozoom   Bildbreite als Vielfaches der Kantenlänge (1 = bündig)
%   \CVphotoshift  Motivversatz; negativ schiebt das Bild nach unten,
%                  das Gesicht rutscht also von oben in die Mitte
\providecommand{\CVphotozoom}{1}
\providecommand{\CVphotoshift}{0}
%   \CVphotomirror 1 = wie aufgenommen, -1 = seitenverkehrt (damit die Person zum
%                  Text hin blickt, wenn das Foto rechts steht)
\providecommand{\CVphotomirror}{1}
\newcommand{\photobox}[2]{%
    \begin{tikzpicture}
        \pgfmathsetlengthmacro{\pbhalf}{0.5*#1}
        \pgfmathsetlengthmacro{\pbwide}{\CVphotozoom*#1}
        \pgfmathsetlengthmacro{\pbshift}{\CVphotoshift*#1}
        \clip[rounded corners=#2] (-\pbhalf,-\pbhalf) rectangle (\pbhalf,\pbhalf);
        \node[anchor=center, inner sep=0] at (0,\pbshift)
            {\scalebox{\CVphotomirror}[1]{\includegraphics[width=\pbwide]{\CVphoto}}};
    \end{tikzpicture}%
}

%-------------------------------------------------------------------------
%  Zweiter Zeitraum derselben Rolle beim selben Arbeitgeber.
%
%  Wer eine Funktion zweimal innehatte, bekommt EINEN Eintrag mit zwei
%  Zeiträumen — nicht zwei Einträge mit identischer Aufgabenliste. Der
%  jüngere Zeitraum steht oben und bestimmt, wo der Eintrag im Lebenslauf
%  einsortiert wird; der ältere folgt als zweite Zeile in der Datumsspalte.
%
%      \cvrole{Immobilien-Bewirtschafterin}{Ing.~Bau~AG}{Vaduz}%
%             {2022 -- 2023\cvalso{2018 -- 2019}}
%
%  Bewusst ohne "und": die Datumsspalte des Akzent-Themes ist 2,2 cm breit,
%  "und 2018 -- 2019" bricht dort auf drei Zeilen um. Dass es zwei Einsätze
%  sind, sagt ohnehin die Zeile unter dem Rollentitel.
%-------------------------------------------------------------------------
\newcommand{\cvalso}[1]{\\ #1}

%-------------------------------------------------------------------------
%  Organisationszeile "Firma · Ort" — beide Teile optional, damit Einträge
%  ohne Arbeitgeber (Familienphase, Auszeit) keinen leeren Trenner zeigen.
%  Nutzt \orgfont, \metafont und die Farbe muted aus dem jeweiligen Theme.
%-------------------------------------------------------------------------
\newcommand{\cvorgline}[2]{%
    % Keine Silbentrennung: Firmen- und Ortsnamen gehören nicht zerlegt.
    % "Zwickau, Deutsch-|land" liest sich wie ein Setzfehler. Umbrüche an
    % den Wortzwischenräumen bleiben möglich.
    \nohyph
    \ifblank{#1}{%
        \ifblank{#2}{}{{\metafont\color{muted}#2}}%
    }{%
        {\orgfont#1}%
        \ifblank{#2}{}{\,{\metafont\color{muted}\textbullet}\,{\metafont\color{muted}#2}}%
    }%
}

%-------------------------------------------------------------------------
%  Links — klickbar, ohne Rahmen
%-------------------------------------------------------------------------
\usepackage[hidelinks, bookmarks=false, unicode]{hyperref}
\urlstyle{same}

%-------------------------------------------------------------------------
%  Weitere optionale Angaben. Leer gelassen erscheinen sie nirgends.
%    \CVkeywords      Stichwortzeile unter dem Titel (kompakt, rueckgrat)
%    \CVwebtext/-url  eigene Website als weiterer Link (rueckgrat)
%    \CVphotoaspect   Höhe des Fotos im Verhältnis zur Breite (rueckgrat),
%                     1 = quadratisch, 1.25 = Hochformat 4:5
%-------------------------------------------------------------------------
\providecommand{\CVkeywords}{}
\providecommand{\CVwebtext}{}
\providecommand{\CVweburl}{}
\newcommand{\cvweb}{\ifdefempty{\CVweburl}{\CVwebtext}{\href{\CVweburl}{\CVwebtext}}}
\providecommand{\CVphotoaspect}{1}

\providecommand{\LblProjects}{Projekte}
\providecommand{\LblReferences}{Referenzen}
\providecommand{\LblEnclosures}{Beilagen}
\providecommand{\LblContents}{Inhalt}
\providecommand{\LblApplicationAs}{Bewerbung als}
\providecommand{\LblCoverLetter}{Anschreiben}
\providecommand{\LblMotivation}{Motivationsschreiben}

%-------------------------------------------------------------------------
%  Beilagenzeile am Ende des Anschreibens. Erscheint nur, wenn die meta.tex
%  \newcommand{\enclosures}{Lebenslauf, Arbeitszeugnisse, Diplome} setzt.
%-------------------------------------------------------------------------
\newcommand{\letterenclosures}{\ifdefined\enclosures
    \par\vspace{0.5cm}\noindent{\metafont\color{muted}\LblEnclosures: \enclosures}\par\fi}

%-------------------------------------------------------------------------
%  Deckblatt (\cvshow{coverpage} in der meta.tex): Briefkopf des Themes,
%  die Zielstelle in der Seitenmitte, unten das Inhaltsverzeichnis des
%  Dossiers. \coverpagecontents setzt application.tex je nach Schaltern.
%-------------------------------------------------------------------------
\providecommand{\coverpagecontents}{\LblCV}
% Grösse der Zielstelle auf dem Deckblatt, vom Theme überschreibbar
\providecommand{\coverpagetitlesize}{\fontsize{24pt}{29pt}\selectfont}
\newcommand{\makecoverpage}{%
    \makeletterhead
    \setlength{\parskip}{0pt}%
    \vspace*{\fill}%
    \noindent{\metafont\color{muted}\MakeUppercase{\LblApplicationAs}}\par\vspace{8pt}%
    \noindent{\nohyph\displayfont\coverpagetitlesize\color{headercolor}\positionname\par}%
    \vspace{6pt}%
    \noindent{\fontsize{13pt}{17pt}\selectfont\companyname\par}%
    \vspace*{\fill}\vspace*{\fill}%
    \noindent{\metafont\color{muted}\MakeUppercase{\LblContents}}\par\vspace{4pt}%
    \noindent{\color{rule}\rule{\linewidth}{0.5pt}}\par\vspace{6pt}%
    \noindent\coverpagecontents\par
    \newpage
}


%-------------------------------------------------------------------------
%  Farben
%-------------------------------------------------------------------------

\definecolor{ink}{HTML}{2A2724}      % Fliesstext
\definecolor{accent}{HTML}{8A4B2A}   % Rost — Sektionen, Daten, Chips
\definecolor{muted}{HTML}{6E6862}    % Sekundärangaben
\definecolor{tint}{HTML}{F4F1EB}     % Sandfläche der Seitenspalte
\definecolor{rule}{HTML}{DED8CE}     % Haarlinien

\colorlet{headercolor}{ink}
\colorlet{linecolor}{accent}

%-------------------------------------------------------------------------
%  Schrift — Source Sans Pro
%-------------------------------------------------------------------------

\setmainfont{SourceSansPro}[
    Extension      = .otf,
    UprightFont    = *-Regular,
    BoldFont       = *-Semibold,
    ItalicFont     = *-RegularIt,
    BoldItalicFont = *-SemiboldIt,
    Ligatures      = TeX,
]
\setsansfont{SourceSansPro}[
    Extension = .otf, UprightFont = *-Regular, BoldFont = *-Semibold,
    ItalicFont = *-RegularIt, Ligatures = TeX]
\newfontfamily\displayfont{SourceSansPro}[
    Extension = .otf, UprightFont = *-Semibold, BoldFont = *-Bold,
    Ligatures = TeX]
\newfontfamily\mediumfont{SourceSansPro}[
    Extension = .otf, UprightFont = *-Semibold, BoldFont = *-Bold,
    Ligatures = TeX]

\makeatletter
\renewcommand\normalsize{\@setfontsize\normalsize{9.4}{13}}
\makeatother
\AtBeginDocument{\normalsize\color{ink}}

\newcommand{\namefont}{\displayfont\fontsize{27pt}{27pt}\selectfont
    \addfontfeature{LetterSpace=-1.5}}
\newcommand{\rolelabelfont}{\fontsize{11pt}{13pt}\selectfont
    \addfontfeature{LetterSpace=8}}
\newcommand{\metafont}{\fontsize{8.4pt}{11.5pt}\selectfont}
\newcommand{\datefont}{\mediumfont\fontsize{8.6pt}{11pt}\selectfont}
\newcommand{\rolefont}{\displayfont\fontsize{11pt}{13.5pt}\selectfont}
\newcommand{\orgfont}{\fontsize{9.4pt}{12pt}\selectfont}
\newcommand{\notefont}{\fontsize{8.8pt}{11.5pt}\selectfont\itshape}
\newcommand{\tagfont}{\fontsize{8pt}{8pt}\selectfont}
\newcommand{\lead}[1]{{\mediumfont#1\cvrealspace}}

\newcommand{\vtopsafe}{\leavevmode}

%-------------------------------------------------------------------------
%  Spaltenraster
%
%  Einziger Stellknopf ist \setsidebarwidth. Inhaltsspalte, Sandfläche und
%  Portraitgrösse leiten sich daraus ab, damit die Bedingung
%
%      Seitenspalte + Steg + Inhaltsspalte = \textwidth
%
%  nicht verletzt werden kann. paracol prüft sie nicht: Eine von Hand
%  falsch gesetzte Spaltenbreite erzeugt keinen Fehler, sondern eine
%  stillschweigend überlaufende Seite.
%
%  Andere Breite wählen — im Vorspann der Bewerbung, nach \input des Themes:
%
%      \setsidebarwidth{4.9cm}
%
%  Untergrenze rund 4,5 cm: In der Kontaktzeile steht die E-Mail-Adresse
%  als nicht trennbarer Block (bei \metafont 111,9 pt) neben der 0,42 cm
%  breiten Icon-Spalte. Darunter ragt sie aus der Seitenspalte heraus.
%  Je schmaler die Spalte, desto zerrissener wird ausserdem das Kurzprofil.
%-------------------------------------------------------------------------

\setlength{\columnsep}{1.5cm}          % Steg, geht in die Rechnung ein

% Die Sandfläche beginnt am Papierrand, nicht am Satzspiegel: Sie läuft
% links in den Anschnitt und endet \bandbleed hinter der Seitenspalte,
% also ungefähr in der Mitte des Stegs. Der linke Seitenrand wird aus der
% geometry-Einstellung berechnet, damit geänderte Ränder nicht von Hand
% nachgezogen werden müssen.
\newlength{\bandbleed}\setlength{\bandbleed}{0.65cm}

% Luft zwischen Portrait und Spaltenkante, je Seite.
\newlength{\photoinset}\setlength{\photoinset}{0.2cm}

\newlength{\sidebarwidth}
\newlength{\contentwidth}
\newlength{\bandwidth}
\newlength{\photodiameter}

\newcommand{\setsidebarwidth}[1]{%
    \setlength{\sidebarwidth}{#1}%
    \setlength{\contentwidth}{\dimexpr\textwidth-\columnsep-\sidebarwidth\relax}%
    \setlength{\bandwidth}{\dimexpr 1in+\hoffset+\oddsidemargin
                                   +\sidebarwidth+\bandbleed\relax}%
    \setlength{\photodiameter}{\dimexpr\sidebarwidth-2\photoinset\relax}%
    \setcolumnwidth{\sidebarwidth, \contentwidth}%
}

\setsidebarwidth{5.6cm}

\columncolor{ink}[0]
\columncolor{ink}[1]

%-------------------------------------------------------------------------
%  Sektionen
%-------------------------------------------------------------------------

% Seitenspalte: kleine Marke mit kurzem Strich darunter
\newcommand{\sbsection}[1]{%
    \par\addvspace{13pt}%
    {\mediumfont\fontsize{8.6pt}{8.6pt}\selectfont
     \addfontfeature{LetterSpace=15}\color{accent}\MakeUppercase{#1}}%
    \par\vspace{4pt}{\color{accent}\rule{1.1cm}{1.1pt}}\par\vspace{6pt}%
}

% Hauptspalte: Marke mit durchgehender Haarlinie
% Optionales Argument wie in kompakt.tex (Platzbedarf in Zeilen;
% hier bleibt die bisherige Voreinstellung von 5 Zeilen.
\newcommand{\cvsection}[2][5]{%
    \par\addvspace{14pt}%
    \Needspace*{#1\baselineskip}% keine Überschrift allein am Seitenfuss
    {\mediumfont\fontsize{9.5pt}{9.5pt}\selectfont
     \addfontfeature{LetterSpace=15}\color{accent}\MakeUppercase{#2}}%
    \par\vspace{4pt}{\color{rule}\rule{\linewidth}{0.6pt}}\par\addvspace{7pt}%
}

% Fakten-Zeile in der Hauptspalte: Marke links, Inhalt rechts
\newcommand{\cvfactwide}[2]{%
    \noindent
    \begin{minipage}[t]{2.4cm}\raggedright\vtopsafe\datefont\color{accent}\strut#1\end{minipage}%
    \hspace{0.4cm}%
    \begin{minipage}[t]{\dimexpr\linewidth-2.8cm\relax}\bodyflow\strut#2\end{minipage}%
    \par\vspace{4pt}%
}

%-------------------------------------------------------------------------
%  Einträge in der Hauptspalte
%-------------------------------------------------------------------------

\newenvironment{cvexperience}{}{\par\addvspace{2pt}}

\newcommand{\cvrole}[4]{%
    \par\addvspace{9pt}\Needspace*{7\baselineskip}\noindent
    \begin{minipage}[t]{\dimexpr\linewidth-3.0cm\relax}%
        \vtopsafe\raggedright\nohyph\rolefont\strut#1\end{minipage}%
    \hfill
    \begin{minipage}[t]{2.9cm}%
        \raggedleft\vtopsafe\datefont\color{accent}\strut#4\end{minipage}%
    \par\vspace{1.5pt}%
    \cvorgline{#2}{#3}%
    % \nobreak vor dem Leim: sonst ist genau diese Stelle -- zwischen
    % Stationskopf und Aufzählung -- ein zulässiger Seitenumbruch, und der
    % Kopf bliebe allein am Fuss der Seite zurück.
    \par\nobreak\vspace{3pt}%
}

\newcommand{\cvnote}[1]{\vspace{-1pt}{\bodyflow\notefont\color{muted}#1\par}\nobreak\vspace{1pt}}

\newcommand{\cvplain}[3]{%
    \par\addvspace{7pt}\bodyflow\noindent
    {\nohyph\rolefont\fontsize{9.8pt}{12pt}\selectfont\strut#1}%
    \,{\metafont\color{muted}\textbullet}\,{\datefont\color{accent}#2}\par
    \vspace{1pt}#3\par%
}

% beginpenalty/midpenalty verbieten den Seitenumbruch vor und innerhalb der
% Liste. Zusammen mit dem \Needspace* in \cvrole bleibt damit jede Station
% als Block zusammen, statt einen einzelnen Punkt auf die Folgeseite zu
% schieben. Preis: Passt eine Station nicht mehr aufs Blatt, wandert sie
% ganz auf die nächste Seite und hinterlässt unten Weissraum.
\newenvironment{cvbullets}{%
    \begin{itemize}[beginpenalty=10000, midpenalty=10000,
                    leftmargin=12pt, labelwidth=6pt, labelsep=6pt, labelindent=0pt,
                    itemindent=0pt, topsep=2.5pt, itemsep=3pt,
                    parsep=0pt, partopsep=0pt, align=left,
                    label={\color{accent}\textbullet\cvrealspace}]
        \bodyflow
    }{\end{itemize}}

%-------------------------------------------------------------------------
%  Fakten-Zeilen in der Seitenspalte — Marke oben, Inhalt darunter
%-------------------------------------------------------------------------

\newcommand{\cvfact}[2]{%
    {\datefont\color{ink}#1}\par\vspace{0.5pt}%
    {\bodyflow\fontsize{8.8pt}{11.8pt}\selectfont\color{muted}#2}%
    \par\vspace{6pt}%
}

%-------------------------------------------------------------------------
%  Schlagwort-Chips
%-------------------------------------------------------------------------

% Höhe und Tiefe des Textfelds werden fest vorgegeben, statt sie vom
% Inhalt bestimmen zu lassen: sonst wird "Excel" (ohne Unterlänge)
% flacher als "Zahlungsverkehr" und die Plaketten sind ungleich hoch.
\newlength{\tagheight}\setlength{\tagheight}{6pt}
\newlength{\tagdepth}\setlength{\tagdepth}{2pt}
\newlength{\tagpad}\setlength{\tagpad}{2.2pt}

\newcommand{\cvtag}[1]{%
    \begin{tikzpicture}[baseline=(tg.base)]
        \node[rectangle, rounded corners=2.5pt, draw=accent!35, fill=accent!7,
              inner xsep=5pt, inner ysep=\tagpad,
              text height=\tagheight, text depth=\tagdepth,
              text=ink, font=\tagfont] (tg) {#1};
    \end{tikzpicture}%
    \hspace{3pt}%
}

% Zeilenabstand der Chip-Zeilen: fest gesetzt, damit die Reihen gleich
% weit auseinanderliegen und nicht von der höchsten Plakette der Zeile
% abhängen. Muss grösser sein als Plakettenhöhe + Tiefe.
\newcommand{\tagblock}[1]{%
    % \raggedright statt des \RaggedRight aus \bodyflow: ragged2e erlaubt am
    % rechten Rand nur 2\,em Dehnung, damit ist ein Umbruch zwischen zwei
    % Plaketten unzulaessig, sobald die Restluecke groesser waere -- TeX laesst
    % die Zeile dann lieber ueberlaufen. \raggedright dehnt unbegrenzt.
    {\bodyflow\raggedright\setlength{\baselineskip}{\dimexpr\tagheight+\tagdepth+2\tagpad+3.6pt\relax}%
     #1\par}%
}

%-------------------------------------------------------------------------
%  Kopf, Kontakt, Portrait
%-------------------------------------------------------------------------

\newcommand{\sidebarphoto}{%
    \noindent\hspace*{\photoinset}\cvstd{photo}{\photobox{\photodiameter}{7pt}}\par\vspace{2pt}%
}

\newcommand{\contactrow}[2]{%
    \noindent\begin{minipage}[t]{0.42cm}\vtopsafe\metafont\color{accent}\strut#1\end{minipage}%
    \begin{minipage}[t]{\dimexpr\linewidth-0.42cm\relax}%
        \vtopsafe\bodyflow\metafont\color{ink}\strut#2\end{minipage}%
    \par\vspace{3.5pt}%
}

\newcommand{\contactblock}{%
    \contactrow{\icoloc}{\CVaddress}%
    \contactrow{\icophone}{\CVphone}%
    \contactrow{\icomail}{\href{mailto:\CVmail}{\CVmail}}%
    \contactrow{\CVlinkicon}{\cvlink}%
    \contactrow{\icocal}{\CVbirth}%
    \contactrow{\icoheart}{\CVcivil}%
    \ifdefempty{\CVnationality}{}{\contactrow{\icoflag}{\CVnationality}}%
    \ifdefempty{\CVpermit}{}{\contactrow{\icopermit}{\CVpermit}}%
    \ifdefempty{\CVgithubtext}{}{\contactrow{\icogithub}{\cvgithub}}%
}

\newcommand{\makenamehead}{%
    {\namefont\nohyph\color{ink}\CVname}\par\vspace{4pt}
    {\rolelabelfont\color{accent}\CVrole}\par\vspace{9pt}
    {\color{rule}\rule{\linewidth}{0.6pt}}\par\vspace{8pt}
    \tagblock{\ContentTags}\vspace{2pt}
}

%-------------------------------------------------------------------------
%  Seitenstil — Sandfläche über die volle Seitenhöhe
%-------------------------------------------------------------------------

\fancypagestyle{cvstyle}{%
    \fancyhf{}
    \renewcommand{\headrulewidth}{0pt}
    \renewcommand{\footrulewidth}{0pt}
    \fancyhead[L]{%
        \begin{tikzpicture}[remember picture, overlay]
            \fill[tint] (current page.north west)
                rectangle ($(current page.south west)+(\bandwidth,0)$);
            % Auf Folgeseiten trägt die Seitenspalte die Kopfzeile
            \ifnum\value{page}>1
                \node[anchor=north west, align=left, inner sep=0]
                    at ($(current page.north west)+(1.3cm,-1.45cm)$)
                    {\mediumfont\fontsize{9.5pt}{12pt}\selectfont\color{ink}\CVname\\[2pt]
                     \metafont\color{muted}\LblCV~~\textbullet~~\LblPage\ \thepage};
            \fi
        \end{tikzpicture}}
    \fancyhead[R]{}
}
\fancypagestyle{letterpage}{\fancyhf{}\renewcommand{\headrulewidth}{0pt}}

%-------------------------------------------------------------------------
%  Briefkopf
%-------------------------------------------------------------------------

\newcommand{\makeletterhead}{%
    \pagestyle{letterpage}\thispagestyle{letterpage}%
    \setlength{\parskip}{7.5pt}%
    \begin{tikzpicture}[remember picture, overlay]
        \fill[tint] ($(current page.north west)$)
            rectangle ($(current page.north east)+(0,-3.5cm)$);
    \end{tikzpicture}%
    \vspace*{0.45cm}
    \noindent
    % 5.6cm ist hier NICHT die Seitenspalte, sondern die Breite, die der
    % rechtsbündige Kontaktblock des Briefkopfs beansprucht. Unabhängig
    % von \sidebarwidth — nicht mitziehen.
    \begin{minipage}[b]{\dimexpr\linewidth-5.6cm\relax}
        {\displayfont\fontsize{21pt}{21pt}\selectfont
         \addfontfeature{LetterSpace=-1}\color{ink}\CVname}\par\vspace{4pt}
        {\fontsize{9.5pt}{11pt}\selectfont
         \addfontfeature{LetterSpace=6}\color{accent}\CVrole}
    \end{minipage}%
    \hfill
    \begin{minipage}[b]{5.4cm}
        \raggedleft\metafont\color{muted}\setlength{\parskip}{0pt}
        \CVaddress\par\CVphone\par\href{mailto:\CVmail}{\CVmail}
    \end{minipage}
    \par\vspace{9pt}
    \noindent{\color{accent}\rule{\linewidth}{1.2pt}}
    \par\vspace{14pt}%
    \letterbodyfont
}

\newcommand{\lettersection}[1]{%
    \par\addvspace{11pt}%
    \noindent{\mediumfont\fontsize{9.5pt}{9.5pt}\selectfont
        \addfontfeature{LetterSpace=13}\color{accent}\MakeUppercase{#1}}%
    \par\vspace{4pt}\noindent{\color{rule}\rule{\linewidth}{0.6pt}}\par\vspace{6pt}%
}

%-------------------------------------------------------------------------
%  Rumpf des Lebenslaufs
%-------------------------------------------------------------------------

\newcommand{\makecvbody}{%
    \pagestyle{cvstyle}\thispagestyle{cvstyle}
    \begin{paracol}{2}

    %--- Seitenspalte -----------------------------------------------------
    \sidebarphoto
    \sbsection{\LblContact}
    \contactblock
    \sbsection{\LblProfile}
    % \par in der Gruppe: sonst endet der Absatz erst nach der Gruppe, und
    % die Trennstrafe aus \bodyflow greift nicht.
    {\bodyflow\fontsize{8.8pt}{12pt}\selectfont\ContentSummary\par}
    \sbsection{\LblLanguages}
    \ContentLanguages

    %--- Hauptspalte ------------------------------------------------------
    \switchcolumn

    \makenamehead

    \cvsection{\LblExperience}
    \ContentExperience

    % Open Source & Projekte (nur, wenn die Inhaltsdatei den Haken füllt)
    {\let\cvfact\cvfactwide\cvprojectsection}

    \cvsection{\LblEducation}
    \ContentEducation

    % Leere Weiterbildung druckt keine Überschrift
    \ifdefempty{\ContentFurther}{}{\cvsection{\LblFurther}\ContentFurther}

    % Kompetenzen ausgeschrieben — im breiten Satzspiegel als Marke/Inhalt-Zeilen
    \cvsection{\LblSkills}
    {\let\cvfact\cvfactwide\ContentSkills}

    {\let\cvfact\cvfactwide\cvextrasections}

    \end{paracol}
}
